\chapter{Conclusiones y Trabajos Futuros}
La metrología monocular es un problema esencial en múltiples ámbitos industriales y
científicos, ya que permite realizar mediciones del entorno físico a partir de una única imagen. Este TFM aborda un método semi-supervisado
capaz de estimar la altura de individuos en escena para entornos no controlados a partir de una sola imagen,
haciendo más accesible el uso de herramientas de metrología para cámaras convencionales. Avances en este
campo podrían extenderse a entornos forenses y de seguridad, la biomecánica o la arquitectura y reconstrucción de escenas~\cite{HoiemPopUp,HoiemObjectsInPerspective,CriminisiApplications, CriminisiReconstruction, CriminisiPaintings, SingleViewExerciseQuantification, ObjectsIntoPhotos}.
\par
En términos prácticos, el objetivo de este TFM ha sido reimplementar y adaptar \emph{ScaleNet}~\cite{SVMIW} ante la
ausencia de una implementación y pesos preentrenados del modelo disponibles públicamente,
empleando el framework~\emph{Detectron2}~\cite{detectron2}. Asimismo, se ha sustituido
el conjunto de datos SUN360~\cite{SUN360}, actualmente no disponible, por Pano360~\cite{SPEC} y se ha diseñado
una pipeline reproducible para entrenamiento y evaluación.
De esta forma, se ha realizado un análisis del método sobre un conjunto de datos
independiente. En este contexto, la principal contribución del trabajo consiste en proporcionar
una implementación reproducible y realizar una evaluación crítica de un método existente de metrología monocular.
\par
En primer lugar, se realizó un estudio de la literatura relativa a la metrología monocular, desde la
estimación del tamaño de objetos en escena a su profundidad dentro de la misma.
Se ha investigado requisitos previos para el uso de métodos clásicos de extracción de información de la escena,
así como métodos que utilizan técnicas de aprendizaje profundo.
Los primeros requieren intervención humana, mediante calibración manual de cámara y conocimiento de la escena
(altura, posición y rotación de la cámara) con respecto a los objetos. Los últimos
aprenden a extraer dicha información de forma automática a partir de los datos.
Durante el análisis del estado del arte, se observa el salto en complejidad teórica al tratar de
utilizar imágenes únicas, con una sola cámara, en lugar de escenas con múltiples puntos de vista
o información adicional de la misma. Ante el bajo número de publicaciones relacionadas,
no se ha identificado una implementación pública disponible que ayude en el avance de la estimación
monocular en entornos no controlados, invitándonos a extender el trabajo de~\cite{SVMIW}.
\par
Para la validación y comparación entre los métodos, se ha utilizado un conjunto de
datos diferente, con diversas escenas de sujetos con altura e información de las cámaras conocidas.
Se hizo uso de diferentes cámaras comunes, así como una cámara equipada con un sensor
de alta calidad para medir la profundidad del sujeto en la escena.
Todas las cámaras fueron cuidadosamente calibradas utilizando el método descrito
en la Sección~\ref{sec:ManualCalibration}.
En total se obtuvo 652 fotogramas válidos con etiquetas reales de la posición de la cámara, del sujeto y su altura.
Entre ellos, hay 12 sujetos diferentes, con alturas entre 1.60 y 1.92 metros.
Se reconoce que, aunque el conjunto incluya un número aparentemente suficiente de ejemplos,
puede observarse una baja diversidad en términos de condiciones ambientales (iluminación y escenarios) y de variabilidad antropométrica
(distribución de morfologías corporales, estatura y complexiones).
Asimismo, las estaturas corresponden a valores autodeclarados por los participantes y no a mediciones realizadas durante
las sesiones de grabación, lo que puede introducir discrepancias entre los valores reportados y las estaturas reales.
En consecuencia, la validación queda condicionada y puede inducir a estimaciones optimistas o pesimistas.
\par
El método clásico hace uso de toda la información disponible de la escena.
Principalmente, la matriz de calibración de la cámara y su orientación, así
como alguna referencia métrica (como puede ser la altura con respecto al suelo).
Utilizando geometría de escena se proyecta el sujeto detectado
en la escena al mundo 3D en escala métrica y se obtiene la estatura.
A diferencia del método clásico, el enfoque basado en aprendizaje profundo no requiere proporcionar
manualmente los parámetros geométricos de la escena durante la inferencia, sino que los estima
implícitamente a partir de la imagen mediante \emph{priors} aprendidos durante el entrenamiento.
Es entrenado en 3 fases, con las dos últimas utilizando 2 conjuntos de datos a
la vez. Uno que permite trabajar sobre la calibración de los parámetros de la escena
(como las rotaciones de la cámara y su campo de visión) y otro que permite aprender
a minimizar el error de reproyección sin saber la altura real de los sujetos o de la cámara en la escena.
\par
Tras el análisis de los resultados, se observa el potencial significativo de este método
de automatizar tareas de metrología. El disponer de una regularización sobre el valor medio
de la altura de los objetos favorece una disminución significativa del error máximo del
método clásico en casos desfavorables (por problemas de calibración o imprecisiones en
las mediciones de referencia de la escena). El error medio del método clásico es de 6.7 centímetros
frente a 9.8 centímetros del método semi-supervisado. El método automático no supera
al clásico en precisión, sino que obtiene un rendimiento en el mismo orden de magnitud sin requerir
calibración manual de la escena. Es decir, el método implementado no
necesita de intervención manual (como realizar la calibración de la cámara o obtener datos
adicionales de la escena) para realizar las estimaciones.
\par
No obstante, se observa que el modelo tiene un efecto de retroceso hacia el~\emph{a priori} de 1.75 metros,
de forma que tiende a realizar subestimaciones para personas de gran estatura y sobreestimaciones para personas de baja estatura.
Además, el error medio de estimación de la altura de la cámara es de 17.9 centímetros.
Este valor es un indicativo de que, aunque el modelo sea capaz de disminuir la pérdida de
reproyección de la escena, los parámetros geométricos estimados no están garantizados de ser
físicamente correctos. Esto podría deberse a que distintas combinaciones de profundidad,
de altura de la cámara y ángulos de rotación de la misma pueden generar un
bajo error de reproyección. Estas son algunas de las limitaciones identificadas del modelo
implementado que abren puertas a posibles futuras investigaciones.
\par
Durante el desarrollo y comprensión de la literatura, se recogió un conjunto de ampliaciones
lógicas adicionales que se pueden realizar sobre el proyecto.
En~\cite{SVMIW} mencionan que una de las principales limitaciones es la necesidad de los objetos de estar situados sobre el suelo.
Además, indican que una posible ampliación sería la estimación de la
región delimitadora del objeto de forma amodal~\cite{kar2015amodalcompletionsizeconstancy}.
Así, se podría extender la aplicación de este sistema a objetos semi-ocultos.
Por otro lado, en~\cite{SAKINA2024765}, se observó que el uso de información adicional como raza y género podría llegar
a mejorar resultados de la estimación. Además, incluir más restricciones junto al error
de reproyección, como utilizando más variables conocidas (por ejemplo utilizando un \emph{a priori} sobre
la altura de la cámara para entrenar un modelo para cámaras exteriores típicas),
podría dar lugar a un modelo con mayor precisión en determinados entornos.
Incluso, se podría optar por utilizar los avances del campo de estimación de profundidad,
como~\cite{yang2024depthv2} capaz de estimar la profundidad en escala métrica,
para estimar nuestra variable $z$ en lugar de intentar estimar la altura de la cámara.
En el campo de estimación de profundidad, hay métodos que utilizan la coherencia entre escenas en un video
(fotogramas vecinos) para entrenar sus modelos, algo que también se podría adaptar a este caso
y disponer así de una regularización más sobre las predicciones.
Por otro lado, existen alternativas que permitan obtener un modelo
más explicable como~\cite{SAKINA2024765}, que utilizan
modelos como regresores lineales y árboles de decisión para estimar la estatura de las
personas, utilizando características extraídas automáticamente de la escena como raza, género
y pose (distancia entre miembros del cuerpo y características faciales).
\par
Con este proyecto se dispone de una implementación del código de~\cite{SVMIW}, con el
caso de estimación de pose para humanos y el caso multiclase (en este caso, como
en el artículo de referencia, se utiliza la clase coche y humano, pero es fácilmente
extensible para cualquier objeto).
Adicionalmente, se realiza una exploración de datos, comparando los métodos de generación
de datos sintéticos de~\cite{SVMIW} y~\cite{SPEC}. Incluyendo la eliminación de imágenes irrelevantes,
gracias al desarrollo de un sistema de visualización y búsqueda de imágenes según características
extraídas automáticamente, permitiendo encontrar imágenes redundantes para reducir el tamaño y ruido del conjunto de datos.
Además, el modelo alcanza la convergencia observada con un menor número de iteraciones bajo la configuración
experimental empleada. Todo ello gracias a optimizaciones varias como aumentación de datos,
tasa de aprendizaje dinámica y por módulo, regularización por recorte del gradiente
y vectorización del proceso de metrología.
\par
Por lo tanto, se concluye que se han completado satisfactoriamente los objetivos planteados, determinando la
posibilidad de resolución del problema y permitiendo una extensión del mismo para diferentes
objetos, conjuntos de datos, arquitecturas de redes convolucionales o métodos de entrenamiento más recientes
(véase el Apéndice~\ref{app:YAMLConfig}),
abriendo así la puerta a futuras investigaciones en el campo. Todo el código disponible está
centralizado en GitHub~\code{https://github.com/briansenas/TFM_SVMIW}, a excepción del conjunto
de datos de validación dado que son datos confidenciales.
